\section*{Overview}
\addcontentsline{toc}{section}{Overview}

\begin{purposebox}
	Seven structured worksheets---one per metalearning step from the background knowledge map.
	Each contains questions that require you to demonstrate understanding by \textbf{explaining in your own words},
	\textbf{working through examples by hand}, and \textbf{connecting concepts to the paper}.

	These are not quizzes. They are thinking tools. The goal is to surface gaps \emph{before} they block you on a milestone.
\end{purposebox}

\subsection*{The Seven Steps}

\begin{table}[h]
	\centering
	\renewcommand{\arraystretch}{1.3}
	\begin{tabularx}{\textwidth}{c l X}
		\toprule
		\textbf{\#} & \textbf{Step}                    & \textbf{Key Deliverable It Supports}           \\
		\midrule
		1           & Mathematical Foundations         & All milestones (prerequisite)                  \\
		2           & Control Theory Fundamentals      & Feb 2: NN Description                          \\
		3           & Stability Analysis               & Feb 9: Existence/Uniqueness; Feb 16: Stability \\
		4           & Neural Network Concepts          & Mar 2: Block Diagram                           \\
		5           & Adversarial Attacks \& Defense   & Apr 6: Simulations (Noise Injection)           \\
		6           & Implementation Skills            & Mar 9: Implementation; Mar 23: Comparison      \\
		7           & Project Timeline \& Deliverables & All milestones (planning)                      \\
		\bottomrule
	\end{tabularx}
\end{table}

\subsection*{How to Use}

\begin{enumerate}[leftmargin=2em]
	\item \textbf{Work through each step in order} (1\,$\to$\,7). Earlier steps build foundations for later ones.
	\item \textbf{Write answers in the provided boxes.}
	\item \textbf{Be honest with the self-assessment checklists}---unchecked items are study targets.
	\item \textbf{Use the ``Notes / Questions'' sections} to capture anything unclear for your next meeting.
	\item \textbf{Revisit as the semester progresses}---your understanding will deepen.
\end{enumerate}

\subsection*{Progress Tracker}

\begin{table}[h]
	\centering
	\renewcommand{\arraystretch}{1.4}
	\begin{tabularx}{\textwidth}{l c c c}
		\toprule
		\textbf{Step}                        & \textbf{Started}  & \textbf{Completed} & \textbf{Confidence (1--5)} \\
		\midrule
		1.\ Mathematical Foundations         & yes               & \rule{2cm}{0.4pt}  & \rule{1.5cm}{0.4pt}        \\
		2.\ Control Theory Fundamentals      & \rule{2cm}{0.4pt} & \rule{2cm}{0.4pt}  & \rule{1.5cm}{0.4pt}        \\
		3.\ Stability Analysis               & \rule{2cm}{0.4pt} & \rule{2cm}{0.4pt}  & \rule{1.5cm}{0.4pt}        \\
		4.\ Neural Network Concepts          & \rule{2cm}{0.4pt} & \rule{2cm}{0.4pt}  & \rule{1.5cm}{0.4pt}        \\
		5.\ Adversarial Attacks \& Defense   & \rule{2cm}{0.4pt} & \rule{2cm}{0.4pt}  & \rule{1.5cm}{0.4pt}        \\
		6.\ Implementation Skills            & \rule{2cm}{0.4pt} & \rule{2cm}{0.4pt}  & \rule{1.5cm}{0.4pt}        \\
		7.\ Project Timeline \& Deliverables & \rule{2cm}{0.4pt} & \rule{2cm}{0.4pt}  & \rule{1.5cm}{0.4pt}        \\
		\bottomrule
	\end{tabularx}
\end{table}
